\documentclass[a4paper,11pt]{article}

\usepackage[a4paper,margin=2cm]{geometry}
\usepackage[T1]{fontenc}
\usepackage[utf8]{inputenc}
\usepackage[ngerman,english]{babel}
\usepackage{lmodern}
\usepackage{microtype}
\usepackage{xcolor}
\usepackage{parskip}
\usepackage{enumitem}
\usepackage[most]{tcolorbox}
\usepackage{titlesec}
\usepackage[hidelinks]{hyperref}

\definecolor{HeaderBlueDark}{RGB}{70,110,170}
\definecolor{RowBlueLight}{RGB}{235,243,255}
\definecolor{RowBlueDark}{RGB}{220,233,250}
\definecolor{SoftGray}{RGB}{90,90,90}

\setlength{\parindent}{0pt}
\setlist[itemize]{leftmargin=1.4em,itemsep=0.55em,topsep=0.4em}
\linespread{1.06}

\titleformat{\section}[block]
{\Large\bfseries\color{HeaderBlueDark}}
{}{0pt}{}
\titlespacing{\section}{0pt}{1.3em}{0.8em}

\newtcolorbox{taskbox}{enhanced,breakable,colback=white,colframe=HeaderBlueDark,boxrule=0.6pt,arc=1.5mm,left=2mm,right=2mm,top=1.5mm,bottom=1.5mm}

\newtcolorbox{infobox}{enhanced,breakable,colback=RowBlueLight,colframe=RowBlueDark,boxrule=0.4pt,arc=1.5mm,left=2mm,right=2mm,top=1.5mm,bottom=1.5mm}

\newcommand{\GermanText}[1]{\textbf{Deutsch: }#1\par}
\newcommand{\EnglishText}[1]{{\small\color{SoftGray}\textbf{English: }#1\par}}

\title{Sollten Menschen häufiger mit öffentlichen Verkehrsmitteln fahren?\\[0.35em]
\large \textnormal{Should people use public transportation more often?}}
\date{}

\begin{document}

   \maketitle

   \section{Thema \textnormal{(Topic)}}

      \begin{taskbox}
         \GermanText{In vielen Städten gibt es immer mehr Autos. Dadurch entstehen Staus, Lärm und Luftverschmutzung. Gleichzeitig investieren viele Städte in Busse, Straßenbahnen und U-Bahnen.}
         \EnglishText{In many cities there are more and more cars. This causes traffic jams, noise, and air pollution. At the same time, many cities are investing in buses, trams, and underground trains.}
      \end{taskbox}

   \section{Aufgabe \textnormal{(Task)}}

      \begin{taskbox}
         \GermanText{Schreiben Sie einen Text über das Thema: \textbf{Sollten Menschen häufiger mit öffentlichen Verkehrsmitteln fahren?}}
         \EnglishText{Write a text about the topic: \textbf{Should people use public transportation more often?}}

         \vspace{0.5em}

         \GermanText{Gehen Sie dabei auf die folgenden Punkte ein:}
         \EnglishText{Address the following points:}

         \begin{itemize}
            \item \textbf{Welche Vorteile haben öffentliche Verkehrsmittel?}\\
            {\small\color{SoftGray}What advantages does public transportation have?}

            \item \textbf{Welche Nachteile haben sie?}\\
            {\small\color{SoftGray}What disadvantages does it have?}

            \item \textbf{Warum benutzen viele Menschen trotzdem lieber das Auto?}\\
            {\small\color{SoftGray}Why do many people nevertheless prefer to use a car?}

            \item \textbf{Was könnte der Staat tun, damit mehr Menschen Bus und Bahn benutzen?}\\
            {\small\color{SoftGray}What could the government do so that more people use buses and trains?}

            \item \textbf{Wie ist Ihre persönliche Meinung?}\\
            {\small\color{SoftGray}What is your personal opinion?}
         \end{itemize}
      \end{taskbox}

   \section{Hinweis \textnormal{(Note)}}

      \begin{infobox}
         \GermanText{Schreiben Sie ungefähr \textbf{150--200 Wörter}. Versuchen Sie, Ihre Meinung klar zu begründen und die Punkte amiteinander zu verbinden.}
         \EnglishText{Write approximately \textbf{150--200 words}. Try to justify your opinion clearly and connect the points with each other.}
      \end{infobox}

      Ich bin persönlich mit den linken Parteien verbunden, deshalb \iffalse Konjunktionaladverbien\fi bin ich ganz für  die Entwicklung der öffentlichen Verkehrsmittel.
      
      (Der erste Vorteil(m))(subject) ist(mit sein, immer Nominativ), dass die einkommensschwachen(kam daramad)  auch eine äquivalente Bewegungsmöglischkeit haben können, damit ich glaube, dass die Regierung(goverment) muss die Reichen mehr besteuren(to tax), damit die Menschen die niedrieger(Komparativ) Einkommen(neutral)haben, können ein effizienter öffentlichen Verkehrsystem erreichen.

      Das bringst viel vorteilen

\end{document}
